\label{sec:introduction}

Every graph-based redistricting algorithm must assign a weight to each edge
in the census-tract adjacency graph before partitioning.
The dominant choice is \emph{geographic boundary length}: the length of the
shared boundary between two adjacent tracts, as computed from TIGER/Line
shapefiles \citep{census2020tiger}.
Longer shared boundaries receive heavier weights; METIS minimizes the total
cut weight, which under geographic weighting corresponds to minimizing total
boundary disruption.

Geographic boundary length is intuitive and well-grounded: long shared
boundaries indicate that two tracts are deeply adjacent, and separating them
imposes a greater perimeter cost on the resulting districts.
But it carries a systematic blind spot.
The weight depends entirely on the geometry of the boundary, not on the
character of the communities on either side.
Two tracts that are perfectly matched economically---both residential suburbs,
both dominated by retail---are weighted identically to two tracts that are
economically opposite---a bedroom suburb next to an industrial logistics hub.
The algorithm cannot distinguish between a natural community boundary and an
artificial cut through homogeneous territory.

This paper proposes a remedy.
We introduce \emph{economic character edge weights}: an additional signal
derived from the Census Bureau's LEHD Workplace Area Characteristics (WAC)
data that encodes the economic profile of each census tract.
Tracts with similar economic profiles receive heavier edges (harder to cut),
while tracts with dissimilar profiles receive lighter edges (easier to cut).
The result is a weight signal that encourages district cuts to run
along natural economic zone transitions---the boundary between a
commercial corridor and the residential neighborhoods surrounding it,
or the line separating an industrial park from adjacent agricultural tracts---
rather than treating all boundaries as fungible.

\paragraph{Motivation from redistricting criteria.}
Most states and redistricting commissions explicitly list
\emph{communities of interest} as a criteria to preserve when drawing
districts \citep{ncsl2021criteria}.
Economic community membership---workers in the same industry sector,
residents sharing a commercial corridor, communities whose livelihoods
depend on the same industrial employer---is one of the oldest and most
legally established forms of community of interest \citep{caprop11,azircprop106}.
Economic character weights operationalize this criterion algorithmically,
using only federal government employment data with no racial or partisan content.

\paragraph{Relationship to existing weight modes.}
The \textsc{bisect} platform supports a three-layer compositor
\citep{bisect2024} in which Layer~2 (edge and vertex weights) is modular.
Economic character is a new \texttt{--weights-override economic-character}
option that composes with the existing geographic, county, and
VRA-aligned weight modes.
It does not change the bisection algorithm (Layer~1) or the seed selection
strategy (Layer~3).

\paragraph{Contributions.}
\begin{enumerate}
  \item An $O(n)$ algorithm for computing per-tract economic character vectors
        from LODES WAC NAICS-sector job counts (Section~\ref{sec:algorithm}).
  \item A cosine-similarity-based edge weight blending formula with well-defined
        behavior for zero-vector (purely residential) tracts
        (Section~\ref{sec:algorithm}).
  \item Empirical evaluation on NC-14, WI-8, and TX-38 demonstrating
        state-specific effects driven by the alignment of economic and
        partisan geographies (Section~\ref{sec:results}).
  \item A legal and practitioner analysis establishing economic character weights
        as a defensible communities-of-interest signal under existing
        redistricting criteria (Section~\ref{sec:legal}).
\end{enumerate}

\paragraph{Paper organization.}
Section~\ref{sec:background} reviews LODES data and related work on economic
geography in redistricting.
Section~\ref{sec:algorithm} defines the algorithm formally.
Section~\ref{sec:results} presents empirical results.
Section~\ref{sec:legal} discusses legal and practitioner considerations.
Section~\ref{sec:conclusion} concludes.
